\section{The Lean formalization}\label{sec:formal}

\subsection{Environment and dependency}
The development lives in \lean{research/grh-pseudorandomness/lean/} of the accompanying repository. It uses Lean v4.35.0 (commit \lean{eb6cafb7}) and
Mathlib commit \lean{56a00585}, the exact pins of the compact edition of the GRH proof,
\lean{github.com/lsernau/grh-lean} at commit \lean{72fead10}, which is a path dependency \cite{grhlean}. All
ten package revisions agree with that edition's manifest. GRH enters in exactly two places:
\begin{itemize}
\item \lean{GRHPseudo.grh}, defined as \lean{LFunction_ne_zero_of_half_lt_re}, of type
  \lean{DirichletLNonvanishingRightOfHalf}
  from the GRH project's \lean{Statements.lean}, which imports only Mathlib;
\item \lean{GRHPseudo.Main.zerosOnLine}, defined as \lean{LFunction_eq_zero_re}, together with
  \lean{riemannHypothesis}.
\end{itemize}
Every other result is proved with these as explicit hypotheses, and the unconditional main theorems are
one-line applications of them. We rebuilt the compact edition from source: 9843 jobs, no errors.

\subsection{Architecture}
\begin{center}\small
\begin{tabular}{p{0.46\textwidth}p{0.46\textwidth}}
\toprule
Module & Content\\
\midrule
\lean{GRH/Interface}, \lean{GRH/Discharge} & hypothesis interface; discharge by the GRH proof\\
\lean{AnalyticNumberTheory/LGrowth} & $\abs{L(s,\chi)}\le1+4q\abs s$ on $\operatorname{Re}s\ge\frac12$\\
\lean{AnalyticNumberTheory/LogDerivGRH} & $L'/L$ bound right of $\frac12$ (uses GRH)\\
\lean{AnalyticNumberTheory/Contour}, \lean{PrimeCharSum} & line shift; Theorem~\ref{thm:A}\\
\lean{Sharp/Twist}, \lean{TwoSided}, \lean{MainTerm} & Lemma~\ref{lem:twosided}; main term from RH\\
\lean{Sharp/GammaBounds}, \lean{FiniteSums} & explicit gamma sides; finite prime sums\\
\lean{Sharp/Assembly} & Theorem~\ref{thm:sharp}\\
\lean{Sharp/KernelBounds}, \lean{ConstBounds}, \lean{Certificate} & Theorem~\ref{thm:closed}\\
\lean{Sharp/Certificate64} & Corollary~\ref{cor:cert}\\
\lean{Sharp/Majorant}, \lean{GammaSharp} & positivity with a majorant; sharper gamma bound\\
\lean{Sharp/ConvKernel}, \lean{ConvAssembly}, \lean{ConvCertificate} & Theorem~\ref{thm:convclosed}, Corollary~\ref{cor:convcert}\\
\lean{Pseudorandomness/*} & bias, Lemma~\ref{lem:eigen}, Lemma~\ref{lem:ranktwo}, Corollary~\ref{cor:A}\\
\lean{MainTheorems/Conditional}, \lean{Unconditional} & the statements of Section~\ref{sec:results}\\
\bottomrule
\end{tabular}
\end{center}
Our own code is about $4400$ lines in $27$ files. A further $2200$ lines are vendored, with attribution, from
\cite{grhmiller,schoenfeldlean,rhprimeerror} and the readable GRH edition.

\subsection{Porting the explicit formula}
The Weil explicit formula of \cite{zeta23} (120 modules, $43{,}000$ lines) targets Lean v4.33.0-rc2 and
Mathlib \lean{51e6992e}. We ported its closure to our pin. Only $12$ files needed changes, in $23$ hunks;
the complete diff against upstream is \lean{verification/zeta23-port.diff}. The changes fall into a few
kinds:
\begin{itemize}
\item a tactic file deleted from Mathlib (\lean{linear_combination'}), which we vendored from the old Mathlib;
\item lemma renames (\lean{Finset.prod_le_prod} $\to$ \lean{prod_le_prod}$_0$, \lean{prod_le_one} $\to$ \lean{prod_le_one}$_0$);
\item \lean{logDeriv_mul}/\lean{logDeriv_prod} applied to pointwise lambdas, which must now be folded first;
\item \lean{positivity} no longer seeing through \lean{id}, and one \lean{fun_prop} measurability goal;
\item \lean{field_simp} normal forms;
\item \lean{autoImplicit} enabled for this library, as upstream.
\end{itemize}
None of the changes alters a statement.

\subsection{Verification}
\begin{center}\small
\begin{tabular}{p{0.58\textwidth}p{0.34\textwidth}}
\toprule
Check & Result\\
\midrule
\lean{#print axioms}, 31 theorems including all of Section~\ref{sec:results} & \lean{propext}, \lean{Classical.choice}, \lean{Quot.sound}\\
all axioms in the environment (Mathlib, GRH proof, port, ours) & Lean core only\\
\lean{leanchecker} replay, 32 project modules & all pass\\
\lean{leanchecker} replay, 120 ported modules & all pass\\
source scan for \lean{sorry}, \lean{axiom}, \lean{native_decide} & none\\
\bottomrule
\end{tabular}
\end{center}
\lean{leanchecker} replays the declarations of each named module; imported modules, including the GRH proof,
are trusted through their \lean{.olean} files, which we built from source. Kernel replay of the GRH proof
itself, including Comparator, nanoda and lean4lean runs, was done by others
\cite{sutherlandgrh,sutherlandshortgrh,rubinsteinpage}. We have not yet run Comparator against a
Mathlib-only statement file for our theorems.

\subsection{Faithfulness of statements}
The definitions that occur in the main statements are short and appear in full in the source:
\lean{bias}, \lean{charSum}, \lean{mass} (\lean{Pseudorandomness/Bias}); \lean{sampleSpace}, \lean{wt},
\lean{cutoff}, \lean{kern}, \lean{testF} (\lean{Sharp/Assembly}); \lean{numB}, \lean{denB}, \lean{UB}, \lean{Kc},
\lean{Cc} (\lean{Sharp/Certificate}, \lean{ConstBounds}); and \lean{y150} (\lean{Sharp/Certificate64}). For the convolution-square results they are \lean{Conv.F}, \lean{Conv.f},
\lean{Conv.K}, \lean{Conv.wt}, \lean{Conv.sampleSpace}, \lean{Conv.cutoff} (\lean{Sharp/ConvKernel}, \lean{ConvAssembly}) and
\lean{Conv.G}, \lean{Conv.UB}, \lean{Conv.numB}, \lean{Conv.denB}, \lean{Conv.y220} (\lean{Sharp/ConvCertificate}). The test
function \lean{testFnN} is grh-miller's and is defined through Mathlib's \lean{Real.smoothTransition}. One
convention differs from the paper: the Lean sample space is $\{0\le n<M\}$, whose term $n=0$ has $\Lambda(0)=0$.
